\documentclass[11pt,a4paper]{article}
\input{../../sablony/pripravy/priprava.sty}
\usepackage{tikz,pgfplots}
\pgfplotsset{compat=1.18}
\nastavPripravu{1. ročník}{Kvadratické rovnice a nerovnice}{124}{Grafy s absolutními hodnotami: kombinace}
\begin{document}
\blok{Co umístit do absolutní hodnoty?}
Někdy je absolutní hodnota kolem \emph{celé kvadratické funkce}; jindy kolem
samotného $x$, kolem lineárního činitele nebo kolem více výrazů současně.
\textbf{Nejde o stejné transformace!} Navážeme na 114. a 115. hodinu.
\blok{Zásady konstrukce grafu}
\begin{enumerate}[label=\arabic*)]
\item Najdi všechny nulové body výrazů, které stojí uvnitř $|\cdot|$.
\item Vyznač na ose intervaly; pro každý urči znaménka a odstraň absolutní hodnoty.
\item Zjednoduš předpis každé větve, sestroj příslušný úsek paraboly či přímky.
\item V hraničních bodech nezapomeň spojit graf (absolutní hodnota je spojitá).
\end{enumerate}
\blok{Společně řešené příklady -- srovnávej se základním grafem}
\begin{prikladReseniBox}[title={Příklad 1 -- modul celého výrazu}]
Porovnejte $g_1(x)=x^2+2x-3$ a $g_2(x)=|x^2+2x-3|$.
\tcblower
\[
\begin{aligned}
g_1(x)&=(x+1)^2-4,\quad Z(g_1)=\{-3;1\},\\
g_2(x)&=\begin{cases}g_1(x),&x\le-3\ \text{nebo}\ x\ge1,\\-g_1(x),&-3<x<1.\end{cases}
\end{aligned}
\]
Část původní paraboly pod osou $x$ se zrcadlí nad ni.
\end{prikladReseniBox}
\begin{prikladReseniBox}[title={Příklad 2 -- modul jen v lineárním členu}]
Sestrojte $g_3(x)=x^2+2|x|-3$ a vysvětlete jeho souměrnost.
\tcblower
\[
\begin{aligned}
g_3(x)&=\begin{cases}x^2-2x-3,&x<0,\\x^2+2x-3,&x\ge0,\end{cases}\\
g_3(-x)&=g_3(x),\quad \text{graf je souměrný podle osy }y.
\end{aligned}
\]
\end{prikladReseniBox}
\begin{prikladReseniBox}[title={Příklad 3 -- dva moduly v jednom předpisu}]
Načrtněte $h_3(x)=|x(x-3)|$ a porovnejte s $h_4(x)=|x|(x-3)$.
\tcblower
\[
\begin{aligned}
h_3(x)&=\begin{cases}x(x-3),&x\le0\ \text{nebo}\ x\ge3,\\-x(x-3),&0<x<3,\end{cases}\\
h_4(x)&=\begin{cases}-x(x-3),&x<0,\\x(x-3),&x\ge0.\end{cases}
\end{aligned}
\]
\textbf{Pozor:} $|AB|=|A|\,|B|$, ale obecně $|AB|\ne|A|B$.
\end{prikladReseniBox}
\blok{Samostatná práce -- vyber podle obtížnosti}
\begin{zadaniReseniBox}[title={Úlohy 1--3}]
\begin{enumerate}[label=\alph*)]
\item $g_4(x)=|x^2+2|x|-3|$.
\item $h_2(x)=x^2-3|x|$.
\item $k_3(x)=|x^2-x|+|x|-2$.
\end{enumerate}
\tcblower
Kritické body pro rozdělení grafu jsou postupně:
\begin{enumerate}[label=\alph*)]
\item $-1,0,1$; funkce je sudá.
\item $0$; funkce je sudá.
\item $0,1$; graf se skládá ze tří kvadratických částí.
\end{enumerate}
Vyznačte v pracovním listu alespoň dva kontrolní body v každém intervalu.
\end{zadaniReseniBox}
\blok{Volitelná rezerva -- skutečně náročný graf}
\begin{zadaniReseniBox}[title={Výzva}]
Sestrojte $k_4(x)=|9-x^2|+|4-x^2|-5$.
\tcblower
\[
\begin{aligned}
k_4(x)&=\begin{cases}2x^2-18,&|x|\ge3,\\0,&2\le|x|\le3,\\8-2x^2,&|x|\le2.\end{cases}
\end{aligned}
\]
Kritické body: $-3,-2,2,3$; střední části jsou vodorovné na ose $x$.
\end{zadaniReseniBox}
\blok{Domácí úkol}
\begin{zadaniReseniBox}[title={Domácí příprava}]
Na přiloženém listu (12 grafů) narýsujte libovolné tři z různých skupin:
\begin{enumerate}[label=\alph*)]
\item $|x^2-3x|$;
\item $x|x-2|+|x^2-2x|$;
\item $|2x+1|-x|x-4|$.
\end{enumerate}
\tcblower
\begin{enumerate}[label=\alph*)]
\item Kritické body $0,3$;
\item kritické body $0,2$ (po úpravě $x|x-2|+|x|\,|x-2|$);
\item kritické body $-\frac12,4$.
\end{enumerate}
Detailní grafy a úplné rozpisy jsou v učitelském klíči pracovního listu.
\end{zadaniReseniBox}
\poznamka{Příklady ze snímku 55 jsou rozděleny na g1--g4, h1--h4, k1--k4. Nemusíš stihnout vše: základem je rozpoznat umístění modulu, ostatní jsou výzva.}
\end{document}
